\chapter{Constructive Scalar Models and OS Data}
\label{ch:constructive-scalar-models-os-data}

\ConventionsEuclidean

Constructive scalar models provide interacting examples in which Euclidean
correlation functions, positivity, growth estimates, and Lorentzian
reconstruction are controlled by theorems in specified dimensions and
parameter regimes.  This chapter records the output data of such a
construction: Schwinger functions, reflection positivity, growth estimates,
locality after reconstruction, and the reconstructed Hilbert-space fields.

\section{Gaussian Reference Field and Wick Powers}

Let \(C=(-\Delta+m^2)^{-1}\) on \(\R^d\) with \(m>0\).  The centered Gaussian
field \(\phi\) is the generalized random field with covariance
\[
  \mathbb E[\phi(f)\phi(g)]
  =
  \int_{\R^d\times\R^d} f(x)C(x-y)g(y)\,d^dx\,d^dy .
\]
\begin{definition}[Wick powers by Gaussian regularization]
\label{def:constructive-wick-powers-gaussian-regularization}
Let \(\phi_\varepsilon=\rho_\varepsilon\ast\phi\) be a smoothed Gaussian
field and let \(C_\varepsilon(0)=\mathbb E[\phi_\varepsilon(x)^2]\), which is
translation independent.  The second and fourth Wick powers at regulator
\(\varepsilon\) are
\[
  :\!\phi_\varepsilon^2\!:
  =
  \phi_\varepsilon^2-C_\varepsilon(0),
\]
\[
  :\!\phi_\varepsilon^4\!:
  =
  \phi_\varepsilon^4
  -6C_\varepsilon(0)\phi_\varepsilon^2
  +3C_\varepsilon(0)^2 .
\]
After smearing against a test function, the Wick power is the \(L^2\)-limit
of these expressions in the Gaussian probability space when the limit exists.
\end{definition}

\section[P(phi) in two dimensions]{\(P(\phi)_2\)}

For a bounded-below polynomial \(P\), the finite-volume interaction in two
Euclidean dimensions is
\[
  V_\Lambda(\phi)
  =
  \int_\Lambda :\!P(\phi(x))\!:\,d^2x .
\]
The finite-volume Schwinger functional is
\[
  S_\Lambda(f_1,\ldots,f_n)
  =
  \frac{1}{Z_\Lambda}
  \mathbb E\!\left[
  \phi(f_1)\cdots\phi(f_n)e^{-V_\Lambda(\phi)}
  \right].
\]
The constructive \(P(\phi)_2\) theorems prove existence of infinite-volume
limits for the Schwinger functions in their stated phase and parameter
regimes, together with Euclidean invariance, symmetry, reflection positivity,
clustering in massive regimes, and the OS growth estimates required for
reconstruction.

The interaction is a Wick-ordered distribution integrated against spacetime
volume.  The infinite-volume and ultraviolet limits are part of the
construction and are controlled before the Schwinger hierarchy is used.

\section[phi-four in three dimensions]{\(\phi^4_3\)}

In three Euclidean dimensions the quartic interaction is
superrenormalizable but requires additional local renormalizations.  A
regularized finite-volume action has the form
\[
  S_{\varepsilon,\Lambda}(\phi)
  =
  \frac{1}{2}\langle\phi,(-\Delta+m^2)\phi\rangle
  +
  \int_\Lambda
  \left[
  \lambda\phi_\varepsilon^4
  +a_\varepsilon\phi_\varepsilon^2
  +b_\varepsilon
  \right]d^3x .
\]
The constants \(a_\varepsilon,b_\varepsilon\) are part of the regulated
Lagrangian.  They are chosen so that Schwinger functions have finite limits
and satisfy the constructive bounds.  The limiting measure is singular with
respect to the original Gaussian reference measure, so the construction is a
limit of regulated measures rather than a fixed-density formula.

\section{OS Data Produced by a Construction}

The output of a constructive scalar theorem is the hierarchy
\[
  S_n(f_1,\ldots,f_n)
  =
  \lim_{\Lambda\uparrow\R^d,\ \varepsilon\downarrow0}
  S_{\varepsilon,\Lambda,n}(f_1,\ldots,f_n)
\]
on Schwartz test functions.  The hierarchy must satisfy:
permutation symmetry; Euclidean covariance; reflection positivity; locality
after reconstruction; cluster or phase-sector hypotheses when a unique
vacuum is intended; and the OS-II regularity and linear-growth assumptions.
Once these are verified, OS reconstruction gives Wightman fields and a
Hilbert space.

The constructive status table in
Chapter~\ref{ch:constructive-status-routes} records examples.  This chapter
supplies the local data format that later lattice, stochastic, and numerical
chapters must match when they claim to approximate a continuum QFT.

\begin{openproblem}[Explicit comparison of constructive routes]
\label{op:constructive-route-comparison}
For a model such as \(\phi^4_3\), compare the Schwinger functions produced by
cluster expansions, stochastic quantization, and rigorous RG constructions in
one common topology.  The comparison should identify the OS Hilbert space,
the local fields, and the parameter map between the constructions.
\end{openproblem}
